\documentclass[nonacm,sigconf]{acmart}

\usepackage{paralist}
\usepackage{xurl}
\usepackage{xspace}

\newcommand{\oursystem}{$\mathsf{ZEE}\tt{200}$\xspace}
\newcommand{\zee}[0]{\ensuremath{\mathsf{ZEE}}\xspace}
\let\P\undefined
\newcommand{\P}{\ensuremath{\mathcal{P}}\xspace}
\newcommand{\V}{\ensuremath{\mathcal{V}}\xspace}

\settopmatter{printacmref=false}
\renewcommand\footnotetextcopyrightpermission[1]{}
\setlength{\emergencystretch}{3em}
\sloppy

\newcommand{\artifactdoi}{\url{https://doi.org/10.5281/zenodo.20575393}}
\newcommand{\artifacttag}{\texttt{ccs2026-ae-v4}}
\newcommand{\artifactcommit}{the commit resolved by \artifacttag}
\newcommand{\passmsg}{\texttt{tight zk cpu test passed}}
\newcommand{\zkvmbin}{\texttt{test\_zkvm\_generic\_asm\_test}}

\begin{document}

\title{ACM CCS '26 Artifact Appendix:\\ \oursystem}
\subtitle{Zero Knowledge for Everything and Everyone @ 200 KHz}

\author{Sunghyeon Jo}
\author{Vladimir Kolesnikov}
\author{Yibin Yang}
\renewcommand{\shortauthors}{Jo et al.}

\maketitle

\appendix

\section{Artifact Appendix}

\subsection{Abstract}

The artifact includes the code and scripts used to reproduce the
evaluation of \oursystem. It contains the \oursystem backend, the reused \zee
baseline, the Docker-based frontend toolchain for compiling {\tt C} programs
to \zee assembly and ZK programs, benchmark inputs, and scripts for
running the prover \P and verifier \V. This appendix explains how to build
the artifact, run the toy test, map the paper's tables to concrete
commands, and interpret the outputs. Successful \oursystem proof runs print
\passmsg. Performance logs report end-to-end time, communication, and
peak resident set size when available. The \oursystem project code is
released under the MIT License. Vendored third-party components remain
subject to their respective licenses.

\subsection{Description \& Requirements}

\subsubsection{Security, privacy, and ethical concerns}

The artifact does not require private credentials and does not
intentionally perform destructive actions. It starts local prover and
verifier processes and opens local network connections. The optional
\texttt{set\_net.sh} script configures Linux traffic control
(\texttt{tc}) and therefore needs administrator privileges. For the main
\oursystem artifact, Docker is needed only for regenerating frontend
artifacts from {\tt C} source. The external SP1 Groth16 reproduction
also invokes Docker through SP1's wrapping toolchain. Docker should be
run with the usual host-security precautions.

\subsubsection{How to access}

The archived artifact submitted for evaluation is:
\begin{compactitem}
\item Zenodo DOI: \artifactdoi
\item GitHub release tag: \artifacttag
\item Git reference: \artifactcommit
\item Repository: \url{https://github.com/ainta/ZEE200}
\end{compactitem}

The DOI above is the Zenodo concept DOI for the artifact. The GitHub
release tag above identifies the exact submitted source snapshot. The
GitHub repository is the maintained development location.

\subsubsection{Hardware dependencies}

The artifact can run on a single machine by running prover and verifier
over localhost. No special hardware is required. The paper's main
measurements use a 2021 ThinkPad X1 Carbon Gen 9 with an Intel Core
i7-1165G7 CPU and 16 GB physical RAM. A comparable Linux machine is
recommended for checking absolute timings. Large \zee baseline
comparisons may need substantially more memory; Table 5 in the paper
uses an AWS \texttt{r6i.8xlarge} instance with 256 GB RAM.

\subsubsection{Software dependencies}

The backend requires CMake 3.10 or newer, \texttt{make}, a {\tt C++17}
compiler, and EMP toolkit dependencies. On an Ubuntu-like system,
\texttt{./setup.sh} installs the expected dependencies, including the
pinned \texttt{emp-tool} and \texttt{emp-ot} dependencies used by the
\zee baseline. Linux is recommended for network experiments because
they use \texttt{tc}. Docker is needed to regenerate frontend assembly
and \texttt{.zk} outputs from the {\tt C} source files. The benchmark
\texttt{.s} and \texttt{.zk} files used in the \oursystem and \zee
experiments are already included in the artifact, so these benchmark
runs do not require the Docker frontend. The external SP1 Groth16
reproduction requires Docker access for the wrapping step.

\subsubsection{Benchmarks}

The artifact includes the workloads reported in the paper:
\texttt{sed} vulnerability payload, Fibonacci with $n=20$ and $n=23$,
\texttt{mergesort} with $n=500$ and $n=5000$, \texttt{gzip} bug
trigger, \texttt{gzip} decompression of a 22 KB compressed image,
\texttt{sha256} with 10 and 200 blocks, and empty-program ZK RAM
sweeps.
Fibonacci $n=20$ and $n=23$ use separate, instruction-count-matched programs.
Their ZEE assembly differs only in the expected-result immediate:
$\operatorname{fib}(20)=6765$ versus $\operatorname{fib}(23)=28657$.
The runner selects the corresponding checked-in artifact from the variant argument.

\subsection{Set-Up}

\subsubsection{Installation}

Download and unpack the Zenodo archive. From the artifact root, build
the \oursystem backend:
\begin{verbatim}
./setup.sh
./scripts/build_release.sh
\end{verbatim}

The release build produces the \oursystem executable under
\texttt{build/bin/}. To build the \zee baseline used by the comparison
scripts:
\begin{verbatim}
cd zee-processor
cmake .
make -j
cd ..
\end{verbatim}

This produces \texttt{zee-processor/ZKMachine}. The \zee scripts use
the checked-in \texttt{.zk} programs under \texttt{benchmarks/}; Docker
is not part of the normal \zee benchmark run path.

\subsubsection{Basic toy test}

Run the prover and verifier in two terminals from the artifact root.
Start the prover first.
\begin{verbatim}
# P terminal
./scripts/run_sed_prover.sh zee200

# V terminal
./scripts/run_sed_verifier.sh zee200
\end{verbatim}

If the build and local networking are correct, both terminals should
print \passmsg. This test exercises assembly parsing, RAM
initialization, prover/verifier networking, the \oursystem proof harness,
and proof verification.

\subsubsection{Expected executable files and scripts}

After installation, the following files are the main evaluator-facing
entry points:
\begin{enumerate}
\item \zkvmbin: the \oursystem \P/\V executable, located in
\texttt{build/bin/}. Direct usage is:
\begin{verbatim}
./build/bin/test_zkvm_generic_asm_test \
  --party PARTY --port 12345 --address ADDR \
  --asm PROGRAM.s --ram-size RAM_BYTES
\end{verbatim}
\texttt{PARTY=1} is the prover and uses \texttt{ADDR=0};
\texttt{PARTY=2} is the verifier and uses the prover address, e.g.,
\texttt{127.0.0.1}.

\item \texttt{zee-processor/ZKMachine}: the \zee baseline binary used
by scripts under \texttt{scripts/zee/}.

\item Top-level wrappers
\texttt{scripts/run\_*\_\{prover,verifier\}.sh}. Each accepts
\texttt{zee200} or \texttt{zee} where the \zee baseline is available.

\item Backend-specific wrappers under \texttt{scripts/zkvm/} and
\texttt{scripts/zee/}. These expose fixed benchmark variants such as
\texttt{run\_zkvm\_fib23\_prover.sh} and
\texttt{run\_zee\_sha200\_verifier.sh}. The top-level
\texttt{run\_fib\_*} and \texttt{run\_sha\_*} wrappers are the preferred
evaluator-facing commands; the prover terminal input selects the
Fibonacci index or SHA iteration count.

\item Frontend helper scripts under \texttt{toolchain/}. These rebuild
\texttt{.s} and \texttt{.zk} files from {\tt C} source through Docker. They
are not needed for running the submitted benchmark files.
\end{enumerate}

\subsection{Evaluation Workflow}

Please run all commands from the artifact root unless otherwise noted.
For a two-party run, execute the prover command on \P first and the
matching verifier command on \V second. In the single-machine setting,
\P listens on localhost and \V connects to \texttt{127.0.0.1}. For
\texttt{fib} and \texttt{sha256}, enter the input values manually in
the prover terminal; do \emph{not} pipe stdin into the prover.

\subsubsection{Reading terminal output}

For \oursystem runs, both parties should print \passmsg. The prover
terminal also prints:
\begin{compactitem}
\item \texttt{[Cleartext] Execution completed: N instructions executed}.
This is the executed extended-\zee instruction count for that benchmark.
\item A final timing line of the form \texttt{T us 1}. Here
\texttt{T} is the prover's end-to-end time in microseconds and may be
printed in scientific notation, e.g., \texttt{1.54163e+06 us 1}.
\item One or more \texttt{communication (B): X} lines. The paper's
communication column reports the protocol communication counter in
bytes converted to MB; use the counter convention used by the artifact
scripts and do not double-count symmetric debug counters if two equal
lines are printed.
\item A peak resident set size line, e.g.,
\texttt{[Linux]Peak resident set size: K}, where \texttt{K} is in KiB
on Linux.
\end{compactitem}

The KHz rate is computed from the raw log values, not from rounded table
entries:
\[
\mathrm{KHz} =
  \frac{N}{T/1{,}000{,}000} \cdot \frac{1}{1000}.
\]
For example, if the prover log reports \texttt{124236 instructions
executed} and \texttt{1.54163e+06 us 1}, the run rate is
$124236 / 1.54163 / 1000 = 80.587$ KHz. Paper table entries are rounded
and averaged over multiple runs, so recomputing from the displayed
seconds and KHz columns will only approximately recover \texttt{N}.

For \zee runs, the verifier prints \texttt{I am convinced!} on success.
The prover log prints \texttt{Num Cycles: C} and \texttt{Num Insts: I}.
Each party then prints a bare seconds value. The \zee baseline does not
print peak RSS itself; wrap prover and verifier commands with
\texttt{/usr/bin/time -v} and parse \texttt{Maximum resident set size
(kbytes): K} when reproducing memory-footprint results.

\subsubsection{Major Claims}
Our artifact makes the following claims.

\textbf{(C1):} \oursystem executes the paper's benchmark suite in ZK and
the verifier accepts successful executions. This is supported by the
benchmark scripts and Table 1.

\textbf{(C2):} \oursystem is substantially faster than \zee on
comparable workloads, with roughly $20$--$40\times$ speedups on the main
examples and about 200 KHz steady-state CPU speed on larger workloads.
This is supported by Table 1 and the larger AWS comparison in Table 5.

\textbf{(C3):} \oursystem has better physical-memory scaling with
configured ZK RAM size than \zee. This is supported by Table 2.

\textbf{(C4):} \oursystem's end-to-end time changes predictably with
bandwidth and latency under the paper's simulated LAN settings. This is
supported by Table 4.

\textbf{(C5):} \oursystem is faster than representative zkVM baselines
on the selected benchmark kernels. This is supported by Table 3. The
submitted artifact directly supports the \oursystem measurements. The
repository also includes pinned reproduction material for the SP1 and
Jolt Table 3 rows under \texttt{zkvm-benchmarks/}.

\subsubsection{Experiments}
We explain how to repeat the experiments that support each claim.

\textbf{(E1):} [Table 1] [10--30 human-minutes + minutes to several
compute hours, depending on selected benchmarks].

\noindent
\underline{Preparation}: build the \oursystem backend. For the paper's
default LAN setting on Linux, configure loopback as:
\begin{verbatim}
IF=lo ./set_net.sh 1000 1
\end{verbatim}
The second argument is one-way delay in milliseconds, so this
approximates 1 Gbps bandwidth and 2 ms RTT. To clear the setting:
\begin{verbatim}
sudo tc qdisc del dev lo root
\end{verbatim}

\noindent
\underline{Execution}: run the following prover/verifier pairs.

\begin{compactitem}
\item \texttt{sed}

\P runs:
\begin{verbatim}
./scripts/run_sed_prover.sh zee200
\end{verbatim}
\V runs:
\begin{verbatim}
./scripts/run_sed_verifier.sh zee200
\end{verbatim}
Paper value: 1.5 s, 82.3 KHz, 53 MB communication.

\item \texttt{fib}, $n=20$

\P runs:
\begin{verbatim}
./scripts/run_fib_prover.sh zee200 20
\end{verbatim}
When prompted in the \P terminal, enter:
\begin{verbatim}
12345
20
\end{verbatim}
\V runs:
\begin{verbatim}
./scripts/run_fib_verifier.sh zee200 20
\end{verbatim}
Paper value: 2.7 s, 172.5 KHz, 136 MB.

\item \texttt{fib}, $n=23$

\P runs:
\begin{verbatim}
./scripts/run_fib_prover.sh zee200 23
\end{verbatim}
When prompted in the \P terminal, enter:
\begin{verbatim}
12345
23
\end{verbatim}
\V runs:
\begin{verbatim}
./scripts/run_fib_verifier.sh zee200 23
\end{verbatim}
Paper value: 8.9 s, 226.3 KHz, 532 MB.

\item \texttt{mergesort}, $n=500$

\P runs:
\begin{verbatim}
./scripts/run_mergesort500_prover.sh zee200
\end{verbatim}
\V runs:
\begin{verbatim}
./scripts/run_mergesort500_verifier.sh zee200
\end{verbatim}
Paper value: 2.6 s, 142.7 KHz, 123 MB.

\item \texttt{mergesort}, $n=5000$

\P runs:
\begin{verbatim}
./scripts/run_mergesort5000_prover.sh zee200
\end{verbatim}
\V runs:
\begin{verbatim}
./scripts/run_mergesort5000_verifier.sh zee200
\end{verbatim}
Paper value: 24.5 s, 199.8 KHz, 1.4 GB.

\item \texttt{gzip} bug trigger

\P runs:
\begin{verbatim}
./scripts/run_gzip_bug_prover.sh zee200
\end{verbatim}
\V runs:
\begin{verbatim}
./scripts/run_gzip_bug_verifier.sh zee200
\end{verbatim}
Paper value: 1.1 s, 16.8 KHz, 21 MB.

\item \texttt{gzip} decompression

\P runs:
\begin{verbatim}
./scripts/run_gzip_decompress_prover.sh zee200
\end{verbatim}
\V runs:
\begin{verbatim}
./scripts/run_gzip_decompress_verifier.sh zee200
\end{verbatim}
Paper value: 17.3 s, 182.1 KHz, 1.0 GB.

\item \texttt{sha256}, 10 iterations

\P runs:
\begin{verbatim}
./scripts/run_sha_prover.sh zee200
\end{verbatim}
When prompted in the \P terminal, enter:
\begin{verbatim}
12345
10
\end{verbatim}
\V runs:
\begin{verbatim}
./scripts/run_sha_verifier.sh zee200
\end{verbatim}
Paper value: 1.6 s, 91.3 KHz, 68 MB.

\item \texttt{sha256}, 200 iterations

\P runs:
\begin{verbatim}
./scripts/run_sha_prover.sh zee200
\end{verbatim}
When prompted in the \P terminal, enter:
\begin{verbatim}
12345
200
\end{verbatim}
\V runs:
\begin{verbatim}
./scripts/run_sha_verifier.sh zee200
\end{verbatim}
Paper value: 19.7 s, 148.7 KHz, 1.1 GB.
\end{compactitem}

\noindent
\underline{Results}: both parties should print \passmsg. The terminal
logs report time and communication. Absolute times depend on machine
load, CPU, memory pressure, and network shaping. The paper reports
averages over 8 runs after discarding the fastest and slowest of 10
total runs.

\textbf{(E2):} [Table 1 and Table 5] [10--30 human-minutes + minutes to
several compute-hours; large cases may require 256 GB RAM].

\noindent
\underline{Preparation}: build \texttt{zee-processor/ZKMachine}.

\noindent
\underline{Execution}: run the same benchmark families with the
\texttt{zee} backend. The \texttt{fib} and \texttt{sha256} prover
commands require manual terminal inputs; enter them in the \P terminal
and do not pipe stdin.

\begin{compactitem}
\item \texttt{sed}

\P runs:
\begin{verbatim}
./scripts/run_sed_prover.sh zee
\end{verbatim}
\V runs:
\begin{verbatim}
./scripts/run_sed_verifier.sh zee
\end{verbatim}
Paper values: Table 1 ZEE 30.1 s; Table 5 AWS ZEE 52.6 s.

\item \texttt{fib}, $n=20$

\P runs:
\begin{verbatim}
./scripts/run_fib_prover.sh zee 20
\end{verbatim}
When prompted in the \P terminal, enter:
\begin{verbatim}
12345
20
\end{verbatim}
\V runs:
\begin{verbatim}
./scripts/run_fib_verifier.sh zee 20
\end{verbatim}
Paper values: Table 1 ZEE 44.5 s; Table 5 AWS ZEE 77.8 s.

\item \texttt{fib}, $n=23$

\P runs:
\begin{verbatim}
./scripts/run_fib_prover.sh zee 23
\end{verbatim}
When prompted in the \P terminal, enter:
\begin{verbatim}
12345
23
\end{verbatim}
\V runs:
\begin{verbatim}
./scripts/run_fib_verifier.sh zee 23
\end{verbatim}
Paper value: Table 5 AWS ZEE 324.6 s.

\item \texttt{mergesort}, $n=500$

\P runs:
\begin{verbatim}
./scripts/run_mergesort500_prover.sh zee
\end{verbatim}
\V runs:
\begin{verbatim}
./scripts/run_mergesort500_verifier.sh zee
\end{verbatim}
Paper values: Table 1 ZEE 43.9 s; Table 5 AWS ZEE 77.3 s.

\item \texttt{mergesort}, $n=5000$

\P runs:
\begin{verbatim}
./scripts/run_mergesort5000_prover.sh zee
\end{verbatim}
\V runs:
\begin{verbatim}
./scripts/run_mergesort5000_verifier.sh zee
\end{verbatim}
Paper value: Table 5 AWS ZEE 1003.1 s.

\item \texttt{gzip} bug trigger

\P runs:
\begin{verbatim}
./scripts/run_gzip_bug_prover.sh zee
\end{verbatim}
\V runs:
\begin{verbatim}
./scripts/run_gzip_bug_verifier.sh zee
\end{verbatim}
Paper values: Table 1 ZEE 5.2 s; Table 5 AWS ZEE 9.15 s.

\item \texttt{gzip} decompression

\P runs:
\begin{verbatim}
./scripts/run_gzip_decompress_prover.sh zee
\end{verbatim}
\V runs:
\begin{verbatim}
./scripts/run_gzip_decompress_verifier.sh zee
\end{verbatim}
Paper value: Table 5 AWS ZEE 1262.1 s.

\item \texttt{sha256}, 10 iterations

\P runs:
\begin{verbatim}
./scripts/run_sha_prover.sh zee
\end{verbatim}
When prompted in the \P terminal, enter:
\begin{verbatim}
12345
10
\end{verbatim}
\V runs:
\begin{verbatim}
./scripts/run_sha_verifier.sh zee
\end{verbatim}
Paper values: Table 1 ZEE 36.1 s; Table 5 AWS ZEE 62.8 s.

\item \texttt{sha256}, 200 iterations

\P runs:
\begin{verbatim}
./scripts/run_sha_prover.sh zee
\end{verbatim}
When prompted in the \P terminal, enter:
\begin{verbatim}
12345
200
\end{verbatim}
\V runs:
\begin{verbatim}
./scripts/run_sha_verifier.sh zee
\end{verbatim}
Paper value: Table 5 AWS ZEE 1210.4 s.
\end{compactitem}

The rows that are dashed in Table 1 are memory intensive on the
16 GB ThinkPad. The paper uses an AWS \texttt{r6i.8xlarge} instance
with 256 GB RAM for the complete Table 5 \zee comparison.

\noindent
\underline{Results}: the comparable Table 1 benchmarks reproduce the
\zee baseline columns: \texttt{sed} 30.1 s, \texttt{fib}
$n=20$ 44.5 s, \texttt{mergesort} $n=500$ 43.9 s, \texttt{gzip} bug
trigger 5.2 s, and \texttt{sha256} 10 iterations 36.1 s. A successful
\zee verifier log contains \texttt{I am convinced!}; a process exit code
alone is not sufficient, because some \zee paths return 0 after printing
\texttt{I am not convinced!}.

\textbf{(E3):} [Table 2] [10--30 human-minutes + minutes to several
compute-hours].

\noindent
\underline{Preparation}: build \oursystem and, if comparing against
\zee, build \texttt{ZKMachine}. For \zee memory measurements, use GNU
\texttt{/usr/bin/time -v} so both parties record peak RSS.

\noindent
\underline{Execution}:
\begin{verbatim}
./scripts/run_empty_ram_sweep_zee200.sh
./scripts/run_empty_ram_sweep_zee.sh
\end{verbatim}

\noindent
\underline{Parsing}: logs are written under \texttt{benchmarks/empty/}.
For \oursystem, parse each RAM section header of the form
\texttt{RAM=BYTES (2\string^e)} in the prover and verifier logs. For
each matching $e$, parse peak RSS from either
\texttt{[Linux]Peak resident set size: K} or GNU time's
\texttt{Maximum resident set size (kbytes): K}. For \zee, parse headers
of the form \texttt{t=T | RAM=2\string^(T+2)} or
\texttt{RAM=2\string^e}; the configured RAM size in Table 2 is
$2^e$ bytes. Parse prover and verifier RSS from GNU time's
\texttt{Maximum resident set size (kbytes): K}.

Table 2 reports total physical memory across both parties:
\[
\mathrm{GiB} =
  \frac{K_{\P} + K_{\V}}{1024^2}.
\]
Treat a row as failing when either party exits nonzero, the success
marker is absent, or the process is killed by memory pressure. In the
paper, \oursystem reaches $2^{27}$ on the 16 GB machine, while \zee is
killed by memory pressure after $2^{24}$.

\textbf{(E4):} [Table 4] [10--30 human-minutes + about 1 compute-hour].

\noindent
\underline{Preparation}: configure Linux loopback for each bandwidth
and one-way delay. For the Table 4 RTT columns, use one-way delays 1 ms,
25 ms, and 50 ms:
\begin{verbatim}
IF=lo ./set_net.sh 100 1
IF=lo ./set_net.sh 500 25
IF=lo ./set_net.sh 1000 50
\end{verbatim}
Change the bandwidth and delay to cover all nine combinations in
Table 4. Clear one setting before installing the next if needed:
\begin{verbatim}
sudo tc qdisc del dev lo root
\end{verbatim}

\noindent
\underline{Execution}: for each setting, run \texttt{mergesort}
$n=500$:
\P runs:
\begin{verbatim}
./scripts/run_mergesort500_prover.sh zee200
\end{verbatim}
\V runs:
\begin{verbatim}
./scripts/run_mergesort500_verifier.sh zee200
\end{verbatim}

\noindent
\underline{Results}: Table 4 reports end-to-end times of 11.3, 13.5,
and 16.2 s at 100 Mbps; 3.5, 5.8, and 9.6 s at 500 Mbps; and 2.6,
5.5, and 9.5 s at 1 Gbps for RTTs 2, 50, and 100 ms respectively.

\textbf{(E5):} [Table 3] [external systems; 10--30 human-minutes +
several compute-hours].

\noindent
\underline{Preparation}: the repository provides pinned reproduction
packages under \texttt{zkvm-benchmarks/}. These are not needed to build
or run \oursystem, but they reproduce the SP1 and Jolt Table 3 baseline
rows. For a clean reproduction, remove stale toolchains or use an empty
\texttt{HOME}, \texttt{CARGO\_HOME}, and \texttt{RUSTUP\_HOME}.

On Ubuntu, install common prerequisites:
\begin{verbatim}
sudo apt-get update
sudo apt-get install -y curl ca-certificates git clang \
  build-essential pkg-config libssl-dev protobuf-compiler \
  gcc-riscv64-unknown-elf
\end{verbatim}

\noindent
\underline{SP1}: the package uses \texttt{succinctlabs/sp1} tag
\texttt{v5.2.4}, locked to commit
\texttt{2a51f3dd370e4c5f74d04dfd89359a13a7e93f99}. Install the SP1
tooling and verify the toolchain:
\begin{verbatim}
# Clean existing SP1 state first if SP1 was already installed:
rm -rf "$HOME/.sp1"
rustup toolchain remove succinct || true

curl -L https://sp1up.succinct.xyz | bash
export PATH="$HOME/.sp1/bin:$PATH"
sp1up -v 5.2.4
cargo +succinct --version
cargo prove --version
docker info
\end{verbatim}

Build the guest ELFs and run the Table 3 proof modes:
\begin{verbatim}
cd zkvm-benchmarks/sp1
unset CARGO_TARGET_DIR
./scripts/build-guests.sh
./scripts/run-stark.sh
./scripts/run-groth16.sh
\end{verbatim}
\texttt{run-stark.sh} runs the core/STARK proof path reported in the
SP1 STARK row; this mode does not provide zero knowledge.
\texttt{run-groth16.sh} runs the Groth16-wrapped proof path reported in
the SP1 Groth16 row and requires Docker. The guest-build script scopes
SP1 guest ELF compiler flags to guest compilation only; do not export
those flags before running the host proof scripts.

\noindent
\underline{Jolt}: the package prepares \texttt{a16z/jolt} tag
\texttt{v0.3.0-alpha}, commit
\texttt{5101ad2143039de6f279613810414c3d071d1f8f}, applies the included
integration patch, and then runs the benchmark cases:
\begin{verbatim}
cd zkvm-benchmarks/jolt
./scripts/prepare-jolt.sh
./scripts/run-all.sh
\end{verbatim}
The first run downloads/registers Jolt's guest toolchain and may
generate a Dory SRS before proving. The included Jolt examples run
Jolt's prove-and-verify path and log \texttt{valid: true} on success;
this default Jolt mode is succinct but does not provide zero knowledge,
and this package does not expose a Groth16 wrapper mode for Jolt.

The expected paper runtimes and deterministic counts are listed in
\texttt{zkvm-benchmarks/expected-results.tsv}.
\begin{verbatim}
cd zkvm-benchmarks
cat expected-results.tsv
\end{verbatim}

\noindent
\underline{Results}: Table 3 reports SP1 and Jolt prover computation
times in seconds. For SP1 STARK, the reported times are 141, 251, 145,
808, 115, and 926 s for fibonacci $n=20$, fibonacci $n=23$,
mergesort $n=500$, mergesort $n=5000$, SHA-256 with 10 iterations, and
SHA-256 with 200 iterations respectively. For SP1 Groth16, the
corresponding times are 2220, 2501, 2212, 3821, 2193, and 3531 s. For
Jolt, the corresponding times are 13, 37, 27, 238, 26, and 195 s.
Absolute timing varies with the machine, but the instruction/cycle
counts and public outputs in \texttt{expected-results.tsv} identify
matching benchmark programs and inputs.

\subsection{Interpreting Deviations}

All benchmark timings are sensitive to CPU frequency scaling, thermal
state, background load, memory pressure, and the exact \texttt{tc}
network configuration. Important checks for AE are:
\begin{compactitem}
\item functional success: \oursystem parties print \passmsg, \zee
verifier prints \texttt{I am convinced!}, SP1 logs verification success,
and Jolt logs \texttt{valid: true};
\item consistency: the same benchmark scripts and inputs correspond to
the paper's benchmark rows;
\item performance trend: longer workloads approach the reported
steady-state speed, \oursystem is substantially faster than \zee on
shared workloads, \oursystem scales to larger RAM settings than \zee,
and lower bandwidth or higher RTT increases runtime.
\end{compactitem}

Known build and run issues are:
\begin{compactitem}
\item The \zee baseline builds as written on a clean Ubuntu
22.04/24.04 environment with the distribution Boost packages. If the
build fails in Boost.Asio code around \texttt{io\_service} or socket
acceptance, check whether a custom Boost installation under
\texttt{/usr/local/include} is shadowing the Ubuntu package.
\item SP1 \texttt{v5.2.4} requires its guest ELF compiler flags for
the RISC-V guest builds, and the provided \texttt{build-guests.sh}
applies them only to guest compilation. If host proof commands fail
with \texttt{OUT\_DIR does not have parent called "target"}, unset
\texttt{CARGO\_TARGET\_DIR} and rerun the SP1 commands.
\item SP1 and Jolt require their pinned external toolchains. If the
corresponding installer or toolchain check fails, rerun the installation
commands in E5 before starting the benchmark scripts.
\item SP1 Groth16 wrapping requires Docker access. If
\texttt{run-groth16.sh} fails with \texttt{Failed to run docker info},
start Docker or fix Docker permissions, then rerun the Groth16 script.
\end{compactitem}

\subsection{Version}

This appendix follows the operational style of the ACM CCS '24
\textit{Tight ZK CPU} artifact appendix while using the ACM CCS 2026
artifact appendix section structure.

\end{document}
